\documentclass[12pt,a4paper]{article}

\usepackage{fontspec}
\setmainfont{Libertinus Serif}
\usepackage[greek.ancient,spanish]{babel}
\babelfont[greek]{rm}{Libertinus Serif}
\usepackage[a4paper,margin=3.2cm]{geometry}
\usepackage{microtype}
\usepackage[hidelinks]{hyperref}

% Griego politónico escrito directamente en UTF-8:
% \gr{...} para palabras sueltas; bloques con \foreignlanguage{greek}{...}
\newcommand{\gr}[1]{\foreignlanguage{greek}{#1}}
\newcommand{\stanzabreak}{\par\vspace{0.6\baselineskip}}

\setlength{\parindent}{1.25em}
\setlength{\parskip}{0.35em}
\raggedbottom


\hypersetup{
  pdftitle={Segundo informe de lectura: Leer la Odisea en tiempos iletrados},
  pdfauthor={}
}

\title{Segundo informe de lectura\\[0.4em]\large\emph{Leer la Odisea en tiempos iletrados}}
\author{}
\date{Agosto de 2026}

\begin{document}

\maketitle

\section*{Dictamen}

El manuscrito funciona ya como un libro, no como una mera recopilación de artículos. Tiene un argumento reconocible, una voz estable y un desenlace preparado desde los primeros capítulos. No necesita una reorganización de fondo ni una adaptación a un lector académico. La revisión pendiente debe concentrarse en la continuidad entre capítulos, la eliminación de repeticiones heredadas de la publicación por entregas y la claridad de algunos pasajes. No conviene neutralizar la voz para ampliar el público.

En su estado actual, el libro consta de veinte capítulos, unas 70.000 palabras y 177 páginas de maqueta. La extensión es proporcionada al recorrido: cine, lectura de Homero, traducción, episodios de la \emph{Odisea}, memoria personal y poemas del autor. Ninguno de esos elementos resulta accesorio al final, aunque no todos tienen el mismo peso en cada capítulo.

\section*{El libro}

La propuesta central consiste en leer la \emph{Odisea} sin aceptar a los dioses como explicación suficiente de la conducta humana. Al retirar esa coartada, Odiseo aparece como un personaje inteligente, valiente, seductor y también responsable de sus decisiones. La película de Nolan proporciona el punto de partida y un adversario concreto; la traducción de Juan Manuel Rodríguez Tobal proporciona el texto sobre el que se discute; la experiencia personal del autor explica por qué esa discusión importa.

Este planteamiento sostiene el conjunto. Los capítulos no se limitan a recorrer aventuras conocidas. La hospitalidad, la guerra, el poder, la violencia contra los débiles, la fidelidad y el remordimiento van modificando el retrato inicial del héroe. La secuencia final ---pretendientes, esclavas, Penélope, Homero y la marcha de Ítaca--- concentra el sentido moral y autobiográfico del libro. El capítulo XX no es un apéndice: completa el argumento al permitir que el Odiseo del autor responda por aquello que el Odiseo homérico no puede reparar.

\section*{Lector y accesibilidad}

El lector natural es culto y curioso, pero no necesita formación clásica. Es, en términos generales, el lector de \emph{Jot Down}: acepta la digresión, reconoce referencias culturales diversas y agradece que la erudición no adopte tono profesoral. El libro explica lo necesario para seguir los episodios y traduce o glosa el griego cuando este interviene en el argumento.

También puede llegar a lectores universitarios. Nolan ofrece una entrada inmediata, pero los asuntos que pueden retenerlos son otros: la fabricación contemporánea de un héroe aceptable, la responsabilidad individual, la violencia de los vencedores y el modo en que una traducción cambia la lectura de un clásico. Las referencias al cine, las series o la cultura popular funcionan porque forman parte del razonamiento. No hace falta añadir reclamos generacionales ni simplificar el vocabulario.

La principal dificultad no es conceptual, sino de densidad. Algunos capítulos combinan muchos nombres, episodios, citas en verso y comentarios antes de recuperar el hilo principal. La solución no es resumir a Homero ni introducir explicaciones escolares, sino comprobar que cada término técnico se aclare en su primera aparición y que el lector sepa en todo momento por qué se cita un fragmento. Las citas extensas deben conservarse cuando el análisis depende de ellas; son parte del libro, no documentación auxiliar.

\section*{La voz}

La voz mezcla lectura, memoria, humor y juicio moral. Ese registro permite hablar de Homero sin fingir distancia académica y enlazar el poema con una vida concreta. Los coloquialismos, las comparaciones inesperadas y las afirmaciones tajantes son parte del contrato con el lector. Suprimirlos de manera sistemática dejaría un ensayo más correcto y menos reconocible.

La accesibilidad debe buscarse en la sintaxis, en la información que se da al lector y en la progresión del argumento, no en rebajar el registro. Tampoco conviene suavizar una afirmación solo porque sea polémica. Sí conviene revisarla cuando se presenta como un hecho y el texto no la sostiene, o cuando su formulación categórica sustituye a un razonamiento que el lector necesita. Esa distinción evita tanto la complacencia como la corrección pedante.

\section*{Puntos de atención}

\begin{enumerate}
  \item \textbf{La serialización.} La adaptación al libro está avanzada, pero todavía hay recordatorios de ideas ya explicadas. No toda recurrencia es un defecto: la responsabilidad de Odiseo, la función de los dioses y el temor del cine a Homero son motivos del libro. Solo deberían eliminarse las recapitulaciones que vuelven a desarrollar un argumento sin añadir una consecuencia nueva.

  \item \textbf{Nolan.} La película da actualidad y establece el contraste inicial. El libro la supera porque acaba tratando asuntos que no dependen de ella. Conviene mantener esa proporción. Cuando Nolan sirve para revelar algo en Homero, cumple una función; cuando reaparece únicamente para reiterar el veredicto, el rendimiento disminuye. El libro debe presentarse como una lectura de la \emph{Odisea}, no como una refutación de la película.

  \item \textbf{La traducción de Tobal.} Es el segundo eje del libro y una de sus singularidades. El comentario de palabras, fórmulas y ritmo permite que el lector perciba decisiones que suelen quedar ocultas en una traducción. No hace falta convertir esos pasajes en filología académica. Sí debe mantenerse una pauta constante: citar, explicar qué hace el verso y regresar al argumento.

  \item \textbf{Los poemas propios.} Su presencia está justificada por el hilo autobiográfico y por la lectura alternativa de Odiseo. El tramo final es donde mejor se integran. Antes de él, cada fragmento debe cumplir una función concreta y no aparecer solo como ilustración de una idea ya expresada en prosa.

  \item \textbf{Las citas.} La abundancia de versos es adecuada para un libro que quiere llevar al lector de vuelta a Homero. La editorial deberá comprobar permisos, extensión y atribución de los fragmentos de traducciones contemporáneas. Es una cuestión práctica, no una objeción literaria.
\end{enumerate}

\section*{Criterio para la edición final}

La última revisión debería aplicar un criterio restrictivo:

\begin{itemize}
  \item corregir errores y aclarar dificultades reales;
  \item retirar una repetición solo cuando no produzca énfasis, ritmo o continuidad;
  \item preservar el humor, los coloquialismos y la primera persona;
  \item no añadir aparato académico que el argumento no necesite;
  \item mantener los versos como parte central de la lectura;
  \item asegurar que los capítulos añadidos enlacen con los anteriores sin volver a presentar el libro.
\end{itemize}

\section*{Conclusión}

El manuscrito tiene forma, lector y razón de ser. Su mejor posibilidad editorial no consiste en hacerlo más neutro, sino en dejarlo más limpio sin borrar su procedencia: un lector de Homero que discute con el poema, con el cine y consigo mismo. Puede resultar accesible a un público joven sin cambiar de voz, siempre que la edición facilite el recorrido y no confunda accesibilidad con simplificación.

La tarea que queda es de ajuste. El libro no necesita otra personalidad ni una tesis más moderada. Necesita que cada capítulo avance y que las reiteraciones propias de las entregas independientes no frenen ese avance.

\end{document}
